\documentclass[border=10pt]{standalone}
\usepackage{tikz,ctex}
\usetikzlibrary{positioning, arrows.meta}

\begin{document}

\begin{tikzpicture}[
    % 全局设置
    node distance = 1.8cm and 1.6cm,
    every node/.style = {
        draw,
        rounded corners,
        thick,
        align = center,
        font = \sffamily\small,
        minimum width = 2.4cm,
        minimum height = 1.0cm
    },
    % 样式定义
    centerStyle/.style = {
        fill = blue!45,
        text = white,
        font = \sffamily\bfseries\large,
        line width = 2pt,
        minimum size = 2.8cm,
        circle
    },
    branchTitle/.style = {
        font = \sffamily\bfseries,
        line width = 1.5pt,
        minimum width = 2.8cm,
        minimum height = 1.0cm
    },
    subNode/.style = {
        fill = white,
        font = \sffamily\small,
        minimum width = 2.0cm,
        minimum height = 0.8cm
    },
    leafNode/.style = {
        fill = white,
        font = \sffamily\footnotesize,
        minimum width = 1.8cm,
        minimum height = 0.7cm
    },
    arrowStyle/.style = {
        ->,
        thick,
        >=stealth,
        shorten >=2pt
    }
]

    % ================= 1. 中心节点 =================
    \node[centerStyle] (center) {可视化训练好的 \\ 卷积网络 \\ Visualizing CNNs};

    % ================= 2. 上方分支：视觉皮层 (红) =================
    \node[branchTitle, fill=red!15, above=of center, xshift=4.0cm, yshift=2.0cm] (cortex) {视觉皮层 \\ Visual Cortex};

    \node[subNode, above left=of cortex, xshift=-1.0cm] (simple) {简单细胞 \\ 边缘/方向};
    \node[subNode, above=of cortex] (gabor) {Gabor 滤波器 \\ 式(10.6)-(10.8)};
    \node[subNode, above right=of cortex, xshift=1.0cm] (complex) {复杂细胞 \\ 池化/不变性};

    \node[leafNode, above=of simple, yshift=0.0cm] (hubel) {Hubel \& Wiesel \\ 1959};
    \node[leafNode, above=of gabor, yshift=0.0cm] (params) {$\theta$ 方向 / $\omega$ 频率 \\ $\phi$ 相位 / $(x_0,y_0)$ 位置};
    \node[leafNode, right=of complex, yshift=0.0cm] (grandmother) {祖母细胞 \\ 深层不变性};
    \node[leafNode, right=of cortex, xshift=0.8cm] (neocognitron) {Neocognitron \\ CNN 前身};

    % ================= 3. 右侧分支：可视化滤波器 (蓝) =================
    \node[branchTitle, fill=blue!15, right=of center] (filters) {可视化滤波器 \\ Visualizing Filters};

    \node[subNode, above right=of filters, yshift=0.3cm] (firstlayer) {第一层可视化 \\ 直接看权重};
    \node[subNode, right=of filters, yshift=1.0cm] (gabor_sim) {与 Gabor 相似 \\ 自然图像统计};
    \node[subNode, right=of filters, yshift=-1.0cm] (maxpatch) {最大激活图像块 \\ 层次结构};
    \node[subNode, below right=of filters, yshift=-0.3cm] (optimize) {数值优化 \\ 生成代表图像};

    \node[leafNode, right=of firstlayer, xshift=0.0cm] (innerprod) {内积幅度大 \\ 激活强};
    \node[leafNode, right=of maxpatch, xshift=0.0cm] (hierarchy) {边缘→纹理 \\ →部件→物体};
    \node[leafNode, right=of optimize, xshift=0.0cm] (regularize) {正则化 \\ 平方和/模糊/裁剪};

    % ================= 4. 下方分支：显著图 (绿) =================
    \node[branchTitle, fill=green!15, below=of center] (saliency) {显著图 \\ Saliency Maps};

    \node[subNode, below left=of saliency, xshift=-0.8cm] (lastconv) {选择最后卷积层 \\ 空间+语义};
    \node[subNode, below=of saliency] (gradcam) {Grad-CAM \\ 式(10.9)-(10.10)};
    \node[subNode, below right=of saliency, xshift=0.8cm] (heatmap) {热力图叠加 \\ 高亮重要区域};

    \node[leafNode, below=of gradcam, yshift=0.0cm] (alpha) {$\alpha_k$ 通道权重 \\ 导数平均};
    \node[leafNode, right=of heatmap, xshift=0.8cm] (example) {狗/猫示例 \\ 图10.15};

    % ================= 5. 左侧分支：对抗攻击 (紫) =================
    \node[branchTitle, fill=purple!15, left=of center] (adversarial) {对抗攻击 \\ Adversarial Attacks};

    \node[subNode, above left=of adversarial, yshift=0.3cm] (definition) {微小修改 \\ 人类不可察觉};
    \node[subNode, left=of adversarial, yshift=0.5cm] (fgsm) {快速梯度符号法 \\ 式(10.11)};
    \node[subNode, left=of adversarial, yshift=-1.5cm] (transfer) {迁移性 \\ 欺骗多个网络};
    \node[subNode, below=of adversarial, yshift=-0.3cm] (physical) {物理对抗 \\ 图10.17};

    \node[leafNode, left=of fgsm, xshift=-0.0cm] (epsilon) {$\epsilon$ 控制幅度 \\ sign(梯度)方向};
    \node[leafNode, left=of transfer, xshift=-0.0cm] (notoverfit) {非过拟合 \\ 线性模型也可};
    \node[leafNode, below left=of physical, yshift=-0.0cm] (security) {安全隐患 \\ 开放研究方向};

    % ================= 6. 左上分支：合成图像 (橙) =================
    \node[branchTitle, fill=orange!15, above=of center, xshift=-5.0cm, yshift=2.0cm] (deepdream) {合成图像 \\ DeepDream};

    \node[leafNode, left=of deepdream, xshift=-3.0cm] (goal) {夸张特征 \\ 放大响应};
    \node[leafNode, above left=of deepdream, xshift=-4.0cm, yshift=-0.0cm] (delta) {$\delta$=预激活 \\ 反向传播};
    \node[leafNode, above left=of deepdream, xshift=0.4cm] (objective) {目标函数 \\ 式(10.12) 平方和};
    \node[leafNode, above=of deepdream, xshift=0.0cm] (art) {艺术生成 \\ 照片/随机噪声};

    % ================= 连线逻辑 =================
    % 中心 -> 各分支标题
    \draw[arrowStyle, red] (center) -- (cortex);
    \draw[arrowStyle, blue] (center) -- (filters);
    \draw[arrowStyle, green!60!black] (center) -- (saliency);
    \draw[arrowStyle, purple] (center) -- (adversarial);
    \draw[arrowStyle, orange] (center) -- (deepdream);

    % 分支标题 -> 子节点 (上方：视觉皮层)
    \draw[arrowStyle, red!60] (cortex) -- (simple);
    \draw[arrowStyle, red!60] (cortex) -- (gabor);
    \draw[arrowStyle, red!60] (cortex) -- (complex);
    \draw[arrowStyle, red!60] (cortex) -- (neocognitron);
    \draw[arrowStyle, red!60] (simple) -- (hubel);
    \draw[arrowStyle, red!60] (gabor) -- (params);
    \draw[arrowStyle, red!60] (complex) -- (grandmother);

    % 分支标题 -> 子节点 (右侧：可视化滤波器)
    \draw[arrowStyle, blue!60] (filters) -- (firstlayer);
    \draw[arrowStyle, blue!60] (filters) -- (gabor_sim);
    \draw[arrowStyle, blue!60] (filters) -- (maxpatch);
    \draw[arrowStyle, blue!60] (filters) -- (optimize);
    \draw[arrowStyle, blue!60] (firstlayer) -- (innerprod);
    \draw[arrowStyle, blue!60] (maxpatch) -- (hierarchy);
    \draw[arrowStyle, blue!60] (optimize) -- (regularize);

    % 分支标题 -> 子节点 (下方：显著图)
    \draw[arrowStyle, green!60!black] (saliency) -- (lastconv);
    \draw[arrowStyle, green!60!black] (saliency) -- (gradcam);
    \draw[arrowStyle, green!60!black] (saliency) -- (heatmap);
    \draw[arrowStyle, green!60!black] (gradcam) -- (alpha);
    \draw[arrowStyle, green!60!black] (heatmap) -- (example);

    % 分支标题 -> 子节点 (左侧：对抗攻击)
    \draw[arrowStyle, purple!60] (adversarial) -- (definition);
    \draw[arrowStyle, purple!60] (adversarial) -- (fgsm);
    \draw[arrowStyle, purple!60] (adversarial) -- (transfer);
    \draw[arrowStyle, purple!60] (adversarial) -- (physical);
    \draw[arrowStyle, purple!60] (fgsm) -- (epsilon);
    \draw[arrowStyle, purple!60] (transfer) -- (notoverfit);
    \draw[arrowStyle, purple!60] (physical) -- (security);

    % 分支标题 -> 子节点 (左上：合成图像)
    \draw[arrowStyle, orange!60] (deepdream) -- (goal);
    \draw[arrowStyle, orange!60] (deepdream) -- (delta);
    \draw[arrowStyle, orange!60] (deepdream) -- (objective);
    \draw[arrowStyle, orange!60] (deepdream) -- (art);

\end{tikzpicture}

\end{document}